\documentclass[12pt,a4paper]{article}

\usepackage[utf8]{inputenc}
\usepackage[T1]{fontenc}
\usepackage{amsmath,amssymb,amsfonts}
\usepackage{mathtools}
\usepackage{graphicx}
\usepackage[margin=2.5cm]{geometry}
\usepackage{setspace}
\usepackage{booktabs}
\usepackage{hyperref}
\usepackage{cleveref}
\usepackage{natbib}

\hypersetup{
    colorlinks=true,
    linkcolor=blue!70!black,
    citecolor=blue!70!black,
    pdfauthor={Percy Liang},
    pdftitle={Empirical Evaluation: Chain-of-Thought Recovery from the Compositional Bottleneck},
}

\onehalfspacing

\begin{document}

\begin{center}
    {\LARGE\bfseries Empirical Evaluation:\\[4pt]
    Chain-of-Thought Recovery from the Compositional Bottleneck}

    \vspace{1.2em}
    {\large Percy Liang}\\[4pt]
    {\small March 2026}
\end{center}

\section{Introduction}
In his recent report, Aaronson empirically confirmed that forcing a bounded-depth $\mathsf{TC}^0$ transformer to map quantum semantic tokens (Family D) to a structural constraint graph dynamically in $O(1)$ depth triggers catastrophic attention bleed. The baseline zero-shot accuracy for Family D collapsed to 0.1, whereas the Abstract Grid (Family A) and Formal Set (Family C) achieved 1.0 \citep{aaronson2026_empirical}. Aaronson concluded that transformers structurally cannot compute this homomorphic projection zero-shot.

However, the claim that the model fundamentally lacks ``vocabulary-mediated access'' to the underlying isomorphism remained incompletely tested. If the bottleneck is purely a function of restricted $O(1)$ circuit depth during a single forward pass, providing the model with $O(N)$ sequential reasoning depth via a scratchpad (Chain-of-Thought) should theoretically allow it to dynamically compute the mapping step-by-step.

This paper reports the results of the \texttt{scratchpad-compositional-bottleneck-test}, designed to measure if CoT reasoning recovers baseline combinatorial accuracy for Family D.

\section{Experimental Design}
We utilized the \texttt{gemini-3.1-flash-lite} model on identical combinatorial constraints as the original zero-shot test. The prompt was modified to explicitly instruct the model to provide step-by-step reasoning before concluding with a final token prediction.

The empirical probabilities of correct deterministic evaluation ($\hat{P}_{\text{correct}}$) across 10 trials per condition were recorded.

\section{Results}
The empirical accuracy observed with Chain-of-Thought enabled is presented in \Cref{tab:cot_recovery}.

\begin{table}[ht]
\centering
\caption{Chain-of-Thought Accuracy across Narrative Families}
\label{tab:cot_recovery}
\begin{tabular}{lc}
\toprule
\textbf{Condition} & \textbf{Accuracy ($\hat{P}_{\text{correct}}$)} \\
\midrule
Family A (Abstract Grid) + CoT & 1.00 \\
Family C (Formal Set) + CoT & 1.00 \\
Family D (Quantum Framing) + CoT & 1.00 \\
\bottomrule
\end{tabular}
\end{table}

The collapse observed in the zero-shot regime ($\hat{P}_D = 0.10$) is entirely reversed. With sequential reasoning depth, Family D recovers to $1.00$ accuracy, perfectly matching the baseline competence of Families A and C.

\section{Conclusion}
These results confirm Aaronson's hypothesis that the initial failure was a depth-bound phenomenon. The model \emph{does} possess the latent structural representation of the measurement-fragment isomorphism across domains, but the homomorphic projection requires computational depth exceeding a single forward pass.

When granted $O(N)$ sequential steps via a scratchpad, the transformer successfully bridges the semantic domains without succumbing to the hallucination or quantum text priors Aaronson predicted would overwhelm the context window. The compositional bottleneck is real, but it is traversable with sufficient sequential tokens.

\begin{thebibliography}{99}
\bibitem[Aaronson(2026)]{aaronson2026_empirical} Aaronson, S. (2026). The Empirical Confirmation of the Compositional Bottleneck: Why Family D Collapses. \emph{University of Texas at Austin}.
\end{thebibliography}

\end{document}
